\documentclass[dvipdfmx, xcolor={dvipsnames},12pt]{beamer}
\usepackage{bxdpx-beamer}
\usepackage{pxjahyper}
\usepackage{otf}

\usepackage[linguistics]{forest}
\usepackage{ascmac}
\usepackage{amsmath}
\usepackage{amsthm}
\usepackage{amssymb}
\usepackage{amsfonts}
\usepackage{latexsym}
\usepackage{mathtools}
\usepackage{float}
\usepackage{url}
\usepackage{xcolor}
\usepackage[most]{tcolorbox}
\usepackage{tikz}
\usepackage{varwidth}
\usetikzlibrary{positioning}

\newtcolorbox{marker}{
enhanced,
colback=red!20,
colframe=red!40!black,
boxrule=0.4pt,
left=8mm,
overlay={
% 左の!帯
\fill[red!80!black] (frame.north west) rectangle ([xshift=6mm]frame.south west);
\node[white,font=\bfseries\Large] at ([xshift=3mm]frame.west) {!};
% 右下折れ曲がり
\fill[red!70!white]
(frame.south east) -- ++(-6mm,0) -- ++(6mm,6mm) -- cycle;
},
drop shadow,
}
\newtcolorbox{dfn}[2][]{%
enhanced,
title={#2},
fonttitle=\bfseries,
colback=cyan!5,
colframe=cyan!50!black,
sharp corners=northwest,
% --- ここからタイトルの見た目調整 ---
attach boxed title to top left={%
xshift=0mm, % 
yshift*=-0.5mm % 
},
varwidth boxed title*=-3mm,
boxed title style={%
colback=cyan!50!black,
colframe=cyan!50!black,
rounded corners,
sharp corners=south,
boxrule=0.5pt,
},
#1
}


\newtcolorbox{thm}[2][]{%
enhanced,
title={#2},
fonttitle=\bfseries,
colback=green!5,
colframe=green!50!black,
sharp corners=northwest,
% --- ここからタイトルの見た目調整 ---
attach boxed title to top left={%
xshift=0mm, % 左端から少し内側
yshift*=-0.5mm % 枠線に半分かぶせる
},
varwidth boxed title*=-3mm, % タイトルの幅に合わせた小さい箱
boxed title style={%
colback=green!50!black,
colframe=green!50!black,
rounded corners,
sharp corners=south,
boxrule=0.5pt,
},
#1
}
\setbeamercolor{block title}{fg=white,bg=blue!70!black}
\setbeamercolor{block body}{fg=black,bg=blue!8}



\title{Chomsky (1956) \S2.4}
\author{荒木 理求}
\date\today

\begin{document}

\begin{frame}
  \titlepage
\end{frame}

\begin{frame}{\S2.4の位置づけ}
ここまで有限状態 Markov 過程では英語を十分には記述できないことを見た. ここで一旦, statistical approximationによるアプローチを検討してみよう. すなわち

\begin{block}{問い}
有限状態 Markov過程をだんだん精密にしていけば, 英語の文法に近づけるのではないか?
\end{block}

つまり一つの完全な文法ではなく,

\[
G_1,\ G_2,\ G_3,\cdots
\]

という近似列を考える.
\end{frame}

\begin{frame}{statistical approximation}
固定した \(n\) に対して, 長さ \(n\) の各列に状態を対応させて文法$G_n$を作る : 

\[
(w_1,\dots,w_n) \longleftrightarrow S_i
\]

ここで \(w_1,\dots,w_n\) は文字ではなく, 英語の語である.

たとえば \(n=2\) なら,

\[
(\text{the},\ \text{man}),\quad
(\text{in},\ \text{the}),\quad
(\text{man},\ \text{is})
\]

のような2語列がある状態に対応するよう, $G_2$を決める :
\[
S_1 \leftrightarrow (\text{the},\ \text{man}),\quad
S_2 \leftrightarrow (\text{in},\ \text{the}),\quad
S_3 \leftrightarrow (\text{man},\ \text{is})
\]
\begin{marker}
\(S_i\) は「直前の \(n\) 語を覚えている状態」と考えられる.
\end{marker}
\end{frame}

\begin{frame}
状態 \(S_i\) が語列$\sigma_i=(w_1,\dots,w_n)$

に対応しているとする.

このとき, 状態 \(S_i\) において次に語 \(X\) を出す確率を
\[
\mathbb{P}(\text{next word}=X \mid \text{state}=S_i)
\coloneq
\mathbb{P}(\text{next word}=X \mid \sigma_i)
\]
と定める.
\begin{marker}
より明示的には
\begin{align*}
&\mathbb{P}(\text{next word}=X \mid \sigma_i)\\
=\;&
\mathbb{P}
\left(
W_{k+1}=X
\mid
(W_{k-n+1},\dots,W_k)=\sigma_i
\right)
\end{align*}
であり, 直前の \(n\) 語を見て, 次の語を確率的に選ぶことを意味する.
\end{marker}
\end{frame}

\begin{frame}
このような文法の出力を
\begin{align*}
&(n+1)\text{オーダーの英語への近似}\\
&\left((n+1)\text{-st order approximation to English}\right)
\end{align*}

と呼ぶ.

\begin{marker}
状態は長さ \(n\) の語列を表し, そこから次の1語を出す.
したがって

\[
\underbrace{w_{k-n+1},\dots,w_k}_{n\text{語}}
\quad + \quad
\underbrace{w_{k+1}}_{1\text{語}}
\]

という合計 \(n+1\) 語の関係を使っている.
\end{marker}
\end{frame}

\begin{frame}
\(n\) を大きくすると, より長い文脈を使って次の語を決めることになる.

\[
\mathbb{P}(W_{k+1} = X \mid w_k)
\]

よりも

\[
\mathbb{P}(W_{k+1} = \mid w_{k-2},w_{k-1},w_k)
\]

の方がより細かい文脈を反映するのは明らかだろう. すなわち\(n\) が大きくなるほど出力はサンプル中の英語に似てくる.

\begin{marker}
ここで得られるのはあくまで「英語らしい語列」であって, 英語の文法を記述したことにはならない.
\end{marker}
\end{frame}

\begin{frame}
\textbf{英語の文法を記述したことにはならない}ことの証左として, Chomskyの有名な例がある :

\begin{enumerate}
\item[(14)] colorless green ideas sleep furiously
\end{enumerate}
これは意味は奇妙だが, 英語として文法的ではある. 一方

\begin{enumerate}
\item[(15)] furiously sleep ideas green colorless
\end{enumerate}
は同じ語から構成されるものの, 非文法的である. \\
いくら近似の精度を上げても, サンプルに依存する限り(14)と(15)を区別することはできない.

\begin{marker}
文法性は単語の出現頻度だけでなく, 語順や構造によって決まる, というのがミソである. $n$-th order approximationは大きい$n$を取ってきたところで, (14)のような正文の多くを排除してしまう.
\end{marker}
\end{frame}

\begin{frame}{有限状態 Markov過程との関係}
\begin{block}{命題}
すべての \(n\)-th order approximation は有限状態 Markov過程として表せるが, 逆は成り立たない.
\end{block}

つまり
\[
\text{\(n\)-th order approximation}
\subsetneq
\text{finite-state Markov process}
\]
である.
\begin{marker}
したがって\(n\)-order approximationの観察だけでは, 文法を記述するシステムとして有限状態Markov過程全体を否定したことにはならない.
\end{marker}
\end{frame}

\begin{frame}
\fbox{証明}\:\:
3つの状態
\[
S_0,\quad S_1,\quad S_2
\]
をもつ有限状態文法を見てみよう. 連接(connected)された状態の組は

\begin{align*}
&(S_0,S_1),\quad (S_1,S_1),\quad (S_1,S_0),\\
&(S_0,S_2),\quad (S_2,S_2),\quad (S_2,S_0)
\end{align*}

であり, それぞれの遷移で

\[
a,\ b,\ a,\ c,\ b,\ c
\]

が出力されるとする.
\end{frame}

\begin{frame}
この過程は次のような列を生成できる : 

\[
a b^2 a,\quad c b^2 c, \quad a b^3 a,\quad c b^3 c, \cdots
\]

しかし
\[
a b^2 c,\quad c b^2 a, \quad a b^3 c,\quad c b^3 a, \cdots
\]

のような列は生成しない. つまり
\[
a\underbrace{b\cdots b}_{m\text{個}}a,
\qquad
c\underbrace{b\cdots b}_{m\text{個}}c
\]

のように, 任意に長い語を挟んだ対応がある. しかし$(n+1)$-th order approximationは$n$語しか見ることができないので, この対応を表せない. \hfill$\square$
\end{frame}

\begin{frame}
\begin{marker}
ここでの dependency は, 前述の\((i, j)\)-dependencyの意味である. 実際, $S \coloneq ab^ma \in L$で
\[
cb^ma \notin L\:\& \: cb^mc \in L
\]
なので, 「$S$は$(1, m+2)$-dependencyを持つ」といえる.
\end{marker}

議論をまとめると : 

\begin{enumerate}
  \item 近似の次数を上げても, 文法的な列と非文法的な列を区別できず, 近似の次数と文法性に相関はない.
  \item \(n\)-th order approximation は finite-state Markov process の特殊な場合にすぎない. したがって, \(n\)-th order approximation への批判は, finite-state Markov process 全体への批判とは区別される.
\end{enumerate}
\end{frame}


\begin{frame}{まとめ}
\S2の議論は, 次のように整理できる :

\begin{enumerate}
  \item 統計的近似は文法性を説明できない.
  \item 有限状態文法は局所的な依存しか扱えない.
  \item 英語には, if--then, either--or などの非局所的依存がある.
  \item それらは入れ子にできるため, mirror-image 型の構造に近い.
  \item 文の長さを有限に制限するのは, 言語理論として本質的でない.
  \item したがって再帰的な文法モデルが必要になる.
\end{enumerate}
次回 : 句構造文法

\end{frame}


\end{document}